\newpage
\pagestyle{plain}
\chapter{Судалгааны арга зүй}

\section{SDLC ба Agile аргачлал}

Судалгааны арга зүй нь боловсролын зорилготой мобайл апп хөгжүүлэлтийн төслийн үндэс болох бөгөөд энэ нь програм хангамжийн хөгжүүлэлтийн амьдралын мөчлөгийн (SDLC - Software Development Life Cycle) загварыг дагаж, Agile аргачлалаар хэрэгжинэ.

Agile нь уян хатан, итератив (давталттай) хөгжүүлэлтэд чиглэдэг бөгөөд хэрэглэгчийн санал хүсэлтэд тулгуурлан байнга шинэчлэх боломж олгодог. Энэ арга зүй нь төслийг үе шатаар хувааж, эрсдлийг бууруулж, чанарыг баталгаажуулдаг.

\subsection{Судалгааны үе шатууд}

\begin{table}[htbp]
\centering
\begin{tabular}{|p{0.8cm}|p{4cm}|p{8cm}|}
\hline
\textbf{№} & \textbf{Үе шат} & \textbf{Тайлбар} \\
\hline
1 & Судалгаа ба шаардлага цуглуулах & SRS баримт бичиг боловсруулах, функционал болон функционал бус шаардлага жагсаах \\
\hline
2 & Дизайн & UML диаграммууд (use case, sequence, ERD) зурах, UI/UX дизайн (Figma) хийх \\
\hline
3 & Хөгжүүлэлт & Flutter ашиглан код бичих, SQLite-д өгөгдөл хадгалах, Firebase-ээр sync хийх \\
\hline
4 & Туршилт & Unit testing, integration testing, usability testing (10-20 хүн) \\
\hline
5 & Үнэлгээ ба шинэчлэлт & KPI хэмжих, A/B testing хийх \\
\hline
6 & Deployment & Google Play, App Store-д байршуулах, CI/CD ашиглах \\
\hline
\end{tabular}
\caption{Судалгааны арга зүйн үе шатууд}
\label{tab:sdlc}
\end{table}

\newpage
\section{Модулийн дизайн}

Аппыг модульчлагдсан архитектурт суурилан хөгжүүлэх бөгөөд энэ нь хөгжүүлэлтийг илүү үр ашигтай болгоно. Модулиудыг бие даасан байдлаар зохион байгуулж, хоорондын хамаарлыг багасгасан.

\begin{table}[htbp]
\centering
\begin{tabular}{|p{3cm}|p{4cm}|p{3cm}|p{3cm}|}
\hline
\textbf{Модуль} & \textbf{Үүрэг} & \textbf{Технологи} & \textbf{Хамаарал} \\
\hline
UI Module & Дэлгэц харуулах & Flutter Widgets & Business Logic \\
\hline
Data Module & Өгөгдөл хадгалах & SQLite & Network \\
\hline
Business Logic & Логик боловсруулах & Dart & UI, Data \\
\hline
Network Module & Онлайн sync & Firebase & Data \\
\hline
\end{tabular}
\caption{Модулийн тохиргоо}
\label{tab:modules}
\end{table}

\section{Өгөгдлийн сангийн дизайн}

Өгөгдлийн сан нь relational байх бөгөөд түүхэн өгөгдлийг хадгалахад зориулагдсан. SQLite ашиглах, entity relationships: One-to-many.

\begin{table}[htbp]
\centering
\begin{tabular}{|p{2cm}|p{7cm}|p{4cm}|}
\hline
\textbf{Хүснэгт} & \textbf{Багана} & \textbf{Тайлбар} \\
\hline
Persons & person\_id (PK), name, birth\_date, death\_date, description, image\_url & Түүхэн хүмүүсийн мэдээлэл \\
\hline
Events & event\_id (PK), title, date, description, person\_id (FK) & Үйл явдлууд \\
\hline
Maps & map\_id (PK), title, coordinates, event\_id (FK) & Газрын зургийн өгөгдөл \\
\hline
Quizzes & quiz\_id (PK), question, answers, correct\_answer & Квиз асуултууд \\
\hline
Media & media\_id (PK), type, url, related\_id (FK) & Мультимедиа файлууд \\
\hline
\end{tabular}
\caption{Өгөгдлийн сангийн хүснэгтийн бүтэц}
\label{tab:database}
\end{table}

\section{Системийн тохиргоо}

\subsection{Front-End}

Flutter framework ашиглан UI/UX дизайн хийх. Интерфейс нь Material Design эсвэл Cupertino загварыг дагаж, Монгол хэл дээр ажиллана. Оффлайн кэшинг ашиглан хурдан ажиллах.

\subsection{Back-End}

Firebase ашиглан өгөгдөл sync хийх. Голчлон оффлайн тул local logic дээр суурилсан, онлайн болоход шинэ контент татах.

\newpage
\subsection{Өгөгдлийн сан}

\begin{itemize}
	\item \textbf{Local:} SQLite
	\item \textbf{Cloud:} Firestore
\end{itemize}

\section{Хамгаалалтын дизайн}

Оффлайн апп-д зориулсан аюулгүй байдлын арга хэмжээ:

\begin{table}[htbp]
\centering
\begin{tabular}{|p{4cm}|p{9cm}|}
\hline
\textbf{Арга хэмжээ} & \textbf{Тайлбар} \\
\hline
AES-256 шифрлэлт & Локал өгөгдлийг шифрлэх \\
\hline
Биометр/PIN & Local authentication ашиглах \\
\hline
Input validation & SQL injection-ээс сэргийлэх \\
\hline
Keystore & API key-г аюулгүй хадгалах \\
\hline
Permissions & Зөвхөн шаардлагатай зөвшөөрөл авах \\
\hline
\end{tabular}
\caption{Хамгаалалтын арга хэмжээнүүд}
\label{tab:security}
\end{table}

\section{UML диаграммууд}

\subsection{Use Case Diagram}


Use case diagram нь апп-ын хэрэглэгчийн үйлдлийг (actor → use case) харуулдаг. Апп-д actor (User) болон use case-ууд (Газрын зураг үзэх, Квиз тоглох, Намтар унших) зурдаг.

\begin{figure}[htbp]
\centering
\includegraphics[width=\textwidth]{use_case.png}
\caption{Use Case диаграмм}
\label{fig:usecase}
\end{figure}
\newpage
\subsection{Sequence Diagram}

Sequence diagram нь UML-ийн хэсэг бөгөөд апп-ын объектуудын хоорондын харилцан үйлдлийг цаг хугацааны дарааллаар харуулдаг. Апп-д хэрэглэгч намтар үзэх үед (User → UI → Provider → DB → Firebase) ямар дарааллаар ажиллахыг тодорхойлдог.

\begin{figure}[htbp]
\centering
\includegraphics[width=\textwidth]{sequence_diagram.png}
\caption{Sequence диаграмм}
\label{fig:sequence}
\end{figure}
\subsection{System Diagram}

System diagram нь апп-ын ерөнхий архитектурыг (компонент, холбоос, өгөгдлийн урсгал) харуулдаг. Апп-д tier-үүдийг (Client: Flutter UI, Data: SQLite + Firebase) тодорхойлдог.

\begin{figure}[htbp]
\centering
\includegraphics[width=0.7\textwidth]{system_diagram.png}
\caption{Системийн диаграмм}
\label{fig:system}
\end{figure}
\newpage
\subsection{ERD Diagram}

ERD (Entity-Relationship Diagram) нь өгөгдлийн сангийн бүтцийг (entity, attribute, relationship) харуулдаг. Persons (name, description) → Events (title, date) one-to-many холбоос зурдаг.

\begin{figure}[htbp]
\centering
\includegraphics[width=\textwidth]{er_diagram.png}
\caption{ERD диаграмм}
\label{fig:erd}
\end{figure}
\section{Хэрэглэгдэх технологи}

\begin{table}[htbp]
\centering
\begin{tabular}{|p{3cm}|p{5cm}|p{5cm}|}
\hline
\textbf{Технологи} & \textbf{Хувилбар} & \textbf{Зориулалт} \\
\hline
Flutter & 3.38+ & Кросс-платформ UI хөгжүүлэлт \\
\hline
Dart & 3.0+ & Програмчлалын хэл \\
\hline
SQLite & 3.x & Локал өгөгдлийн сан \\
\hline
Firebase & Latest & Cloud sync, authentication \\
\hline
Figma & Latest & UI/UX дизайн \\
\hline
Git/GitHub & Latest & Version control, CI/CD \\
\hline
\end{tabular}
\caption{Хэрэглэгдэх технологийн жагсаалт}
\label{tab:tech}
\end{table}